% Generated draft: requires mathematical and source review.
\documentclass[11pt]{article}
\usepackage[margin=25mm]{geometry}
\usepackage{fontspec}
\setmainfont{Latin Modern Roman}
\usepackage{amsmath,amssymb,mathtools,unicode-math}
\setmathfont{Latin Modern Math}
\usepackage{enumitem,needspace,xurl,hyperref}
\setlength{\parindent}{0pt}
\setlength{\parskip}{5pt}
\setlength{\emergencystretch}{3em}
\newcommand{\N}{\mathbb{N}}
\newcommand{\Z}{\mathbb{Z}}
\newcommand{\Q}{\mathbb{Q}}
\newcommand{\R}{\mathbb{R}}
\newcommand{\C}{\mathbb{C}}
\newcommand{\F}{\mathbb{F}}
\newcommand{\A}{\mathbb{A}}
\newcommand{\PP}{\mathbb{P}}
\newcommand{\E}{\mathbb{E}}
\newcommand{\eps}{\varepsilon}
\newcommand{\dd}{\,\mathrm{d}}
\begin{document}
\section*{UnsolvedMath ID: 30002483}
\textbf{Source problem code:} catalogue entry
\subsection{Definitions and mathematical statement}
Let \(G\) be quasisplit, let \(B\subseteq G\) be a Borel subgroup, and let \(U\) be its ``nilradical,'' in the terminology of the supplied record. Let \(\psi\) be a nondegenerate character of \(U\), let \(H\subseteq G\) be reductive and spherical, and put \(X=G/H\). The action under consideration is the left action on \(X\). Here spherical conventionally means that a Borel subgroup has an open orbit on the homogeneous space after passage to an algebraic closure; the record does not specify the base field or its convention for spherical spaces. To interpret characters and group orbits, \(U\) would ordinarily be the unipotent radical of \(B\); a Lie algebra nilradical requires a different formulation. Neither interpretation is explicitly resolved by the supplied statement.

The source asks whether
\[
\sum_{x\in O} S^*(Ux)^{U,\psi}
\quad\text{is dense in}\quad
S^*(X)^{U,\psi},
\]
and describes the summands as one-dimensional spaces. Superscript \((U,\psi)\) denotes distributions transforming by the character \(\psi\) under the induced \(U\)-action. The record does not define \(S^*\), its equivariance convention, or the topology used for density. In an Archimedean interpretation, a possible convention is the continuous dual of an appropriate Schwartz test-function space, but this is not specified by the source record.

There is a material mismatch in the orbit parameter: the record defines \(O\) to be the union of the open \(B\)-orbits in \(G\), although the target distribution space is on \(X\). An orbit \(Ux\) for \(x\in G\) is an orbit in \(G\), and the record supplies no map placing its distributions in \(S^*(X)\). Consequently the displayed density question is not fully defined as written.

A conditional formulation of the apparent intended question is as follows. If \(O\) is intended to be
\[
O_X=\bigcup\{\mathcal O\subseteq X:\mathcal O\text{ is an open }B\text{-orbit}\},
\]
fix the underlying field and point spaces, the meaning of openness, a distribution theory \(S^*\), a topology \(\tau\) on \(S^*(X)^{U,\psi}\), and well-defined equivariance-preserving maps
\[
\iota_x:S^*(Ux)^{U,\psi}\longrightarrow S^*(X)^{U,\psi}
\qquad(x\in O_X).
\]
For every datum satisfying the stated hypotheses and these specified conventions, is
\[
\overline{\operatorname{span}\!\left(
\bigcup_{x\in O_X}\iota_x\bigl(S^*(Ux)^{U,\psi}\bigr)
\right)}^{\,\tau}
=S^*(X)^{U,\psi}?
\]
The span consists of finite linear combinations. This conditional formulation does not establish that replacing \(G\) by \(X\) is the intended correction, that the maps \(\iota_x\) exist, or that all the orbit spaces are one-dimensional. These points require source clarification, including any stabilizer compatibility required for nonzero \((U,\psi)\)-equivariant orbit distributions.
\subsection{Short English statement}
The proposed question asks whether Whittaker-equivariant distributions on the spherical quotient \(G/H\) can be approximated by finite sums of distributions arising from \(U\)-orbits in the open \(B\)-orbits. The supplied wording instead places those open orbits in \(G\), and leaves the field, distribution spaces, topology, and orbit-to-quotient maps unspecified. The precise displayed reformulation is therefore conditional. The supplied status is partially solved; this is a dated curation label, not an independently verified mathematical conclusion.
\subsection{Sources}
\begin{itemize}
\item Local UnsolvedMath registry: \texttt{data/unsolvedmath/problems.json}; supplementary registry: \texttt{catalog/problems.json}. Supplied permanent ID: \(30002483\). Exact source problem code: \texttt{catalogue entry}. Title: \emph{Density of Whittaker Distributions on Spherical Spaces}. Collection: Oberwolfach Reports Open Problems. No individual proposer is identified in the supplied record.
\item Supplied report citation: \emph{Representation Theory and Analysis of Reductive Groups: Spherical Spaces and Hecke Algebras} (2014), Oberwolfach Reports, pp.\ 145--212, DOI 10.4171/owr/2014/03. \url{https://doi.org/10.4171/owr/2014/03}. The supplied background describes these reports as documenting workshops and their open-problem discussions; no page specifically locating this problem is supplied.
\item Supplied literature source: Aizenbud--Gourevitch, spherical-space distributions. \url{https://arxiv.org/abs/1311.6111}.
\item Supplied literature source: Aizenbud--Gourevitch--Sayag--Kemarsky, z-finite distributions. \url{https://arxiv.org/abs/1405.2540}.
\item Supplied literature source: Oberwolfach representation theory report. \url{https://publications.mfo.de/bitstream/handle/mfo/3393/OWR_2014_03.pdf?isAllowed=y&sequence=1}.
\item Literature assessment dated exactly \texttt{2026-08-21}: the supplied curation labels the problem partially solved and reports vanishing results away from open \(B\)-orbits and related density results under additional hypotheses, while describing the full Archimedean spherical-space density conjecture as unresolved. This is supplied dated curation, not independently verified. The record itself requests independent verification of time-sensitive status and cited claims.
\item The supplied statement status is \texttt{corrected\_verified}, accompanied by a claim of manual semantic, typographic, and LaTeX review. That registry assertion does not resolve the formulation issues identified here or constitute mathematical verification of this draft.
\end{itemize}
\end{document}
